
\subsubsection{Research Questions}

We decompose the core hypothesis into four testable questions:

\textbf{Q1 (h-e1): Mechanism Functionality.} Does the tri-modal framework correctly implement predicted weight patterns? We validate that weight scheduling operates as designed, feedback collectors function, and aggregation produces measurable rewards.

\textbf{Q2 (h-m1): Phase 1 Execution Foundation.} Does execution-heavy weighting in Phase 1 (0--30\%) drive fastest correctness improvement? We predict execution weight highest among signals, pass@1 improvement rate exceeds later phases by $\geq 8\times$, and execution weight correlates negatively with progress.

\textbf{Q3 (h-m2): Phase 2 Scalable Quality.} Does AI feedback peak in Phase 2 (30--70\%) enable quality refinement without correctness regression? We predict AI weight peaks around 50\% and exceeds other signals, quality improves $\geq 0.05$, and pass@1 maintains $\geq 95$\% of Phase 1 endpoint.

\textbf{Q4 (h-m3): Phase 3 Edge Case Tuning.} Does increasing human feedback weight in Phase 3 (70--100\%) prevent execution-only collapse? We analyze conflict cases (execution succeeds but quality low) and predict preference scores resolve to [0.1, 0.4] rather than collapsing below 0.1.

\subsubsection{Dataset}

We use HumanEval (164 hand-written programming problems from OpenAI) and MBPP (874 crowd-sourced Python programming problems from Google Research), totaling 1,038 competitive programming tasks. Both datasets provide automated test cases for execution feedback. We use fixed random seed 42 for reproducibility.

\subsubsection{Validation Approach}

Gate criteria focus on mechanism verification rather than absolute performance. For h-e1, we require code runs without errors, mechanism implements correctly, and metrics are measurable. For h-m1/m2/m3, we define quantitative thresholds on weight dominance patterns, improvement rates, and correlation coefficients. These gates test whether predicted patterns emerge.

We compare tri-modal dynamic scheduling against single-feedback baselines (execution-only, AI-only, human-only) as conceptual comparisons. However, since we do not perform actual training, all models use the same pretrained checkpoint. This establishes the validation framework is operational rather than providing performance comparison.
